\section{Finite relational witnesses and backward viability}
\label{sec:relational}

\subsection{The coherence problem}

For each precision \(p\), a proof may produce at least one valid finite witness,
but witnesses at different precisions may be incompatible.  Choosing the least
candidate independently at each level does not solve the problem.  It may jump
between branches and fail to define a name.  Conversely, appealing to an
infinite path through a finitely branching tree would import a choice or
compactness principle stronger than the finite theorem being proved.

The remedy is to select only a finite path to the finite horizon currently
requested.

\subsection{Finite tower}

Fix a horizon \(h\).  For each level \(p\le h\), let
\[
 C_p=\{a<B_p:\Cand(p,a)=1\}
\]
be a finite candidate cloud.  Let \(E_p(a,b)\) be the compatibility relation
from level \(p\) to \(p+1\).

Define backward viability relative to \(h\):
\begin{align*}
 V_h(h,a)&:\equiv \Cand(h,a),\\
 V_h(p,a)&:\equiv \Cand(p,a)\wedge
   \exists b<B_{p+1}\,[E_p(a,b)\wedge V_h(p+1,b)].
\end{align*}
All existential quantifiers are bounded and are implemented by Boolean folds.

\begin{theorem}[Viability/path equivalence]
\label{thm:viability}
For \(p\le h\), \(V_h(p,a)\) holds if and only if there is a compatible path
\[
 a=a_p,a_{p+1},\ldots,a_h
\]
with \(a_j\in C_j\) and \(E_j(a_j,a_{j+1})\) for every \(j<h\).
\end{theorem}

\begin{proof}
Use backward induction on \(h-p\).  At \(p=h\), viability is exactly candidate
validity and the one-element path suffices.  For the induction step, unfold
\(V_h(p,a)\).  A viable successor \(b\) exists; by the induction hypothesis it
extends to a compatible tail, which we prepend by \(a\).  Conversely, the
second element of any compatible full path is a viable successor by the
induction hypothesis, so the defining bounded existential holds.
\end{proof}

\subsection{Least viable path}

Define \(a_p\) to be the least viable candidate at the initial level.  Given
\(a_j\), choose the least \(a_{j+1}<B_{j+1}\) satisfying
\[
 E_j(a_j,a_{j+1})\wedge V_h(j+1,a_{j+1}).
\]
By \cref{thm:viability}, such a candidate exists whenever the initial
viability test succeeds.  Bounded least search is primitive recursive.

\begin{corollary}[Finite dependent choice]
For every finite horizon, a nonempty viable tower has a primitive-recursively
selected compatible path.  This construction uses only bounded search and
quantifier-free induction; it does not imply an infinite choice principle.
\end{corollary}

\subsection{Stable quotient erasure}

Often the theorem does not need the internal candidate; it needs an observable
\(W_h(a_0,\ldots,a_h)\).  If every viable path has the same observable, then
choice is irrelevant.

\begin{theorem}[Exact stable erasure]
Assume that for any two viable paths \(\alpha,\beta\) to horizon \(h\),
\[
       W_h(\alpha)=W_h(\beta).
\]
Then the output computed from the least viable path is independent of every
tie-breaking convention and is extensionally characterized by the finite
relation alone.
\end{theorem}

\begin{proof}
The least viable path is viable.  Apply the stability assumption to it and to
an arbitrary viable path.  Their observables are equal.  Since all quantifiers
range over bounded path codes, the argument Booleanizes into a finite proof.
\end{proof}

Exact equality can be replaced by a metric diameter bound:
\[
 d(W_h(\alpha),W_h(\beta))\le\varepsilon_h.
\]
Then every tie-breaking convention yields outputs within \(\varepsilon_h\).
This is the natural statement for approximate roots, eigenvalues, local
solutions, and geometric selections.

\subsection{Changing horizons}

The path selected for horizon \(h+1\) need not extend the path selected for
horizon \(h\).  This does not matter if the level-\(p\) candidates lying on
full viable paths to either horizon have small diameter.  Define the selected
observable at precision \(p\) using a horizon \(H(p)\ge p\).  Suppose:
\begin{enumerate}
\item full-horizon viable candidates at level \(p\) have diameter at most
\(\delta_p\);
\item every compatibility edge changes the represented value by at most
\(\mu_p\);
\item a primitive-recursive tail bound controls
\(\sum_{j\ge p}\mu_j\).
\end{enumerate}
Then selections made with different future horizons are close.  The next
section turns this fact into a regular name.

\subsection{Applications}

\paragraph{Polynomial roots.}
At precision \(p\), candidates are finite rational disks or root clusters.
Edges express refinement and preservation of multiplicity.  At collisions,
there is no ordered label; symmetric observables or cluster data are stable.
A separation certificate permits exact labels locally.

\paragraph{Local charts.}
Candidates are chart indices and approximate coordinates.  Edges assert that
overlap transition maps agree within a rational budget.  A downstream theorem
uses a stable observable, not a globally canonical chart.

\paragraph{Finite sheaf data.}
Candidates are local sections satisfying equations on a finite cover.  Edges
are restriction compatibility.  A gluing certificate may return a cloud of
global representatives; quotient equality erases nonuniqueness.

\paragraph{Mordell--Weil descent.}
Candidates can encode rational points on covering spaces together with maps to
the base curve.  The framework verifies compatibility and descent.  It does
not manufacture rational representatives from abstract Selmer classes; that
remains a producer obligation.

\subsection{Lean plan}

The file \code{RelationalWitnesses.lean} defines finite bounds, validity,
edges, and viability at the semantic level and records object-language targets.
A complete migration should:
\begin{enumerate}
\item encode the level/horizon dynamic program as fixed \(\PR\) symbols;
\item prove the viability recursion equations;
\item Booleanize path validity;
\item prove both directions of \cref{thm:viability} by quantifier-free
induction on the remaining horizon;
\item define least viable successor with the shared bounded-search library;
\item formulate stability as a local finite theorem and compose it with the
selected path proof.
\end{enumerate}
